\subsection{可逆模类群}

*（我们保持第 5 节与第 6 节的记号。）*

在乘法之下，B 的 A-子模组成一个交换幺半群 M，以 A 为单位元。于是可逆模就是 M 的可逆元，从而组成一个交换群 J。我们已看到（第 6 节命题 10），\(M \in J\) 的逆是 A: M。

设 A*（相应地 B*）是 A（相应地 B）的可逆元乘法群，并用 u 表示典范单射 A → B。对每个 \(b \in B^{*}\) ，\(\theta(b) = bA\) 是一个可逆 A-子模。映射 \(\theta: B^{*} \to J\) 是一个同态，其核为 \(u(A^{*})\) ；其余核记作 C 或 C(A)。群 C 称为 B 的可逆 A-子模类群。我们构造出了下面的正合列

\[
(1) \longrightarrow A ^ {*} \stackrel {u} {\longrightarrow} B ^ {*} \stackrel {\theta} {\longrightarrow} \mathfrak {J} \stackrel {\rho} {\longrightarrow} \mathfrak {C} \longrightarrow (1)\tag{10}
\]

其中 (1) 表示仅由单位元素组成的群，而 \(\rho\) 是典范映射 \(\mathfrak{J} \to \mathfrak{C} = \mathfrak{J}/\theta(\mathrm{B}^*)\) 。

由于 B 的每个可逆 A-子模 M 都是秩为 1 的投射模（第 6 节定理 4），元素 \(\operatorname{cl}(M) \in P(A)\) 有定义（第 4 节）。

命题 12. 映射 cl: \(\mathfrak{J} \to P(A)\) 通过取商定义了一个从 \(\mathfrak{C} = \mathfrak{J}/\theta(B^{*})\) 到典范同态

\[
\phi \colon P (\mathrm{A}) \to P (\mathrm{B})
\]

（第 4 节）的核上的同构。

换言之，有正合列

\[
(1) \longrightarrow \mathrm{A} ^ {*} \stackrel {u} {\longrightarrow} \mathrm{B} ^ {*} \stackrel {\theta} {\longrightarrow} \mathfrak {J} \stackrel {\mathrm{cl}} {\longrightarrow} P (\mathrm{A}) \stackrel {\phi} {\longrightarrow} P (\mathrm{B}).\tag{11}
\]

由第 6 节命题 11 以及 \(\mathbf{P}(\mathbf{A})\) 中加法的定义可得，关系式 \(\operatorname{cl}(\mathbf{M},\mathbf{N}) = \operatorname{cl}(\mathbf{M}) + \operatorname{cl}(\mathbf{N})\) （对 \(\mathbf{M}\) 、\(\mathbf{N}\) ，均在 \(\mathfrak{J}\) 中）成立，这表明 \(\operatorname{cl}\) 是一个同态。若 \(\mathbf{M} \in \mathfrak{J}\) 同构于 \(\mathbf{A}\) ，则存在 \(b \in \mathbf{B}\) 使 \(\mathbf{M} = Ab\) ，而且由于 \(\mathbf{M}\) 是可逆的，存在 \(b' \in B\) 使 \(b'b = 1\) ，换言之 \(b\) 在 \(\mathbf{B}\) 中可逆；其逆是直接的。因此 \(\operatorname{cl}\) 在 \(\mathfrak{J}\) 中的核是 \(\theta(\mathbf{B}^*)\) 。

现在来确定 cl 的像。若 \(M \in J\) ，则 \(M \otimes_{A} B = S^{-1}M = B\) （第 5 节命题 8 (c)），因此 \(\operatorname{cl}(M) \in \operatorname{Ker}(\phi)\) 。反之，设 P 是秩为 1 的投射 A-模，使得 \(P_{(B)} = P \otimes_{A} B\) B-同构于 B。由于 P 是平坦 A-模，单射 \(u: A \to B\) 定义了一个单射 \(u \otimes 1: P \to P_{(B)} = B\) ，于是 P 等同于 B 的一个 A-子模；依据第 5 节命题 8 (c)，P 是非退化的，而第 6 节定理 4 表明 P 是可逆的。因此 \(\phi\) 的核等于 \(\operatorname{cl}: J \to P(A)\) 的像。

推论 1. B 的两个可逆 A-子模在 C 中有相同的像，其充分必要条件是它们同构。

推论 2. 若环 B 是半局部的，则 B 的可逆 A-子模类群 C 典范地等同于秩为 1 的投射 A-模类群 P(A)。

这时 \(\mathrm{P(B)} = 0\) （第 3 节命题 5）。

注. 推论 2 的假设在下面两种情形下成立：

\begin{enumerate}
\def\labelenumi{(\arabic{enumi})}
\tightlist
\item
  A 是整环，S 是 A 中由 \(\neq0\) 的元素组成的集合，这时 B 是 A 的分式域。在这种情形，B 的可逆 A-子模也称为可逆分式理想；其中作为单生成自由 A-模 Ab（ \(b\neq0\) ，b 在 B 中）的那些，恰好就是《代数》第 VI 章 §1 第 5 节中定义的分式主理想。
\end{enumerate}

*(2) 环 A 是 Noether 环，而 S 是 A 中不是零因子的元素组成的集合，使得 B 是 A 的全分式环。这时 \(S = A - \bigcup_{i} p_i\) ，其中诸 \(p_i\) 是 Ass(A) 的元素（个数有限）（第 IV 章 §1），因此 B 是半局部的（§3 第 5 节命题 17）。*


